\documentclass[12pt]{article}

% =========================
% 0) Metadata (EDIT ME)
% =========================
\newcommand{\CourseCode}{MATH 521}
\newcommand{\CourseTitle}{Analysis I}
\newcommand{\CourseSection}{LEC 002}
\newcommand{\Semester}{Spring 2026}
\newcommand{\Instructor}{Thierry Laurens}
\newcommand{\MeetingInfo}{MWF 1:20--2:10 \;|\; VV B119} % optional
\newcommand{\HomeworkNumber}{10} % <-- change each week
\newcommand{\YourName}{Alejandro De La Torre}
\newcommand{\DueDate}{April 10, 2026} % <-- or set e.g. January 31, 2026

% =========================
% 1) Core math + proof tools
% =========================
\usepackage{amsmath, amssymb, amsthm}

% =========================
% 2) Lists and spacing
% =========================
\usepackage{enumitem}
\setlist[enumerate,1]{label=\textbf{(\alph*)}, leftmargin=2em, itemsep=0.25em}
\setlist[itemize]{leftmargin=1.5em, itemsep=0.25em}
\setlength{\parskip}{0.35em}
\setlength{\parindent}{0pt}

% =========================
% 3) Page geometry + header/footer
% =========================
\usepackage{geometry}
\geometry{margin=1in}

\usepackage{fancyhdr}
\pagestyle{fancy}
\fancyhf{}
\lhead{\CourseCode\ --- \CourseSection, \Semester \;}
\rhead{Homework \HomeworkNumber}
\cfoot{\thepage}
\renewcommand{\headrulewidth}{0.4pt}
% Fixes common fancyhdr warning about header height:
\setlength{\headheight}{15pt}
% Optional: reclaim a bit of vertical space after increasing headheight
\addtolength{\topmargin}{-2pt}

% =========================
% 4) Links (subtle)
% =========================
\usepackage[hidelinks]{hyperref}
\setcounter{tocdepth}{1}

% =========================
% 5) Quotes / figures
% =========================
\usepackage{csquotes}
\usepackage{graphicx}
\graphicspath{{images/}} % put figures in ./images/

% =========================
% 6) Simple scratch box (no tcolorbox)
% =========================
% A lightweight environment for brief planning / scratch work.
% This avoids any tcolorbox dependencies.
\newenvironment{scratch}{%
  \par\smallskip
  \noindent\textbf{Discussion \& Scratch Work}\begin{quote}\small
}{%
  \end{quote}\smallskip
}

% =========================
% 7) Lightweight theorem-like environments (optional)
% =========================
\newtheorem*{claim}{Claim}
\newtheorem*{lemma}{Lemma}
\newtheorem*{remark}{Remark}
\newtheorem{theorem}{Theorem}
\newtheorem{proposition}{Proposition}
\newtheorem{corollary}{Corollary}
\newtheorem{definition}{Definition}
\newtheorem{example}{Example}
\newtheorem{exercise}{Exercise}

% =========================
% 8) Problem header command (with TOC + PDF bookmarks)
% Usage: \problem[Optional short title]{<number>}
% =========================
\makeatletter
\newcommand{\problem}[2][]{%
  \def\prob@title{Problem #2\if\relax\detokenize{#1}\relax\else\ (\textit{#1})\fi}%
  \phantomsection%
  \pdfbookmark[1]{\prob@title}{prob#2}%
  \addcontentsline{toc}{section}{\prob@title}%
  \section*{\prob@title}%
  \vspace{-0.25em}\hrule\vspace{0.8em}%
}
\makeatother

% =========================
% 9) Title block
% =========================
\title{Homework \HomeworkNumber}
\author{\YourName \\
\CourseCode\ --- \CourseTitle \\
\CourseSection \\
Instructor: \Instructor}
\date{\DueDate}

\begin{document}
\maketitle

% ------------------ collaboration + sources ------------------
% \section*{Collaboration Statement}
% Write “I worked alone” or list collaborators / external resources.
% "I worked on this homework independently. No outside collaborators or resources were used."

\tableofcontents
\bigskip
\hrule
\bigskip

\newpage

% ============================================================
% Problems (duplicate blocks as needed)
% ============================================================

\problem[]{1}
Let $(X,d)$ be a metric space and let $Y\subseteq X$. We call the set $Y$
complete if the subspace $Y$ with the induced metric is complete. Prove that if
$Y$ is complete, then $Y$ is closed in $X$. (Hint: It suffices to show
$\overline{Y}\subseteq Y$. Use sequences.)

\begin{scratch}
The key definitions here come straight from lecture. First, saying that $Y\subseteq X$ is \emph{complete} means that every Cauchy sequence in the subspace $(Y,d_Y)$ converges to some point of $Y$, where $d_Y$ is the induced metric. Since the induced metric is just the restriction of $d$ to pairs of points in $Y$, a sequence in $Y$ is Cauchy in $(Y,d_Y)$ exactly when it is Cauchy in $(X,d)$.

The most important theorem from lecture for this problem is the sequential characterization of closure: for a subset $A\subseteq X$,
\[
x\in \overline{A}
\quad\Longleftrightarrow\quad
\text{there exists a sequence } \{a_n\}\subseteq A \text{ such that } a_n\to x.
\]
We also proved that every convergent sequence in a metric space is Cauchy, and that limits in a metric space are unique.

So the strategy is very direct. Since the hint says it suffices to prove $\overline{Y}\subseteq Y$, I start with an arbitrary point $x\in\overline{Y}$. By the closure theorem, there is a sequence $\{y_n\}\subseteq Y$ with $y_n\to x$ in $X$. Because convergent sequences are Cauchy, $\{y_n\}$ is Cauchy. Since $Y$ is complete, that same sequence must converge to some point $y\in Y$. But now the sequence $\{y_n\}$ converges to both $x$ and $y$ in the metric space $X$, so uniqueness of limits forces $x=y$. Hence $x\in Y$. Since every point of $\overline{Y}$ lies in $Y$, we get $\overline{Y}\subseteq Y$, and therefore $Y$ is closed.
\end{scratch}

\textbf{Proof.}\\
We will show that
\[
\overline{Y}\subseteq Y.
\]
Since always $Y\subseteq \overline{Y}$, this will imply $Y=\overline{Y}$, and hence $Y$ is closed in $X$.

So let $x\in \overline{Y}$. By the sequential characterization of closure proved in lecture, there exists a sequence $\{y_n\}_{n\ge 1}$ in $Y$ such that
\[
y_n\to x
\qquad \text{in } (X,d).
\]
Since every convergent sequence in a metric space is Cauchy, the sequence $\{y_n\}$ is Cauchy in $(X,d)$. Because each $y_n$ lies in $Y$ and the induced metric on $Y$ is just the restriction of $d$, it follows that $\{y_n\}$ is also a Cauchy sequence in the subspace $(Y,d_Y)$.

Now $Y$ is complete by assumption, so every Cauchy sequence in $(Y,d_Y)$ converges to a point of $Y$. Therefore there exists some $y\in Y$ such that
\[
y_n\to y
\qquad \text{in } (Y,d_Y).
\]
But convergence in the subspace metric means exactly the same distance condition as convergence in $(X,d)$, so in fact
\[
y_n\to y
\qquad \text{in } (X,d).
\]
Thus the sequence $\{y_n\}$ converges in $(X,d)$ to both $x$ and $y$. By uniqueness of limits in metric spaces, we must have
\[
x=y.
\]
Since $y\in Y$, it follows that $x\in Y$.

Because $x\in \overline{Y}$ was arbitrary, we conclude that
\[
\overline{Y}\subseteq Y.
\]
Therefore $Y$ is closed in $X$.
\qed

\newpage

\problem[]{2}
Let $(X,d)$ be a metric space. Prove that if $A,B\subseteq X$ are both compact
then $A\cup B$ is compact.

\begin{scratch}
The main definition from lecture that matters here is the definition of compactness: a subset $K\subseteq X$ is compact if every open cover of $K$ has a finite subcover. So for this problem, I think I should start with an arbitrary open cover of $A\cup B$ and try to show that only finitely many of those open sets are needed.

The big thing is that any open cover of $A\cup B$ automatically gives an open cover of $A$ and also an open cover of $B$, since both sets are contained in $A\cup B$. Because $A$ is compact, I can extract finitely many sets from the cover that already cover $A$. Likewise, because $B$ is compact, I can extract finitely many sets from the same cover that cover $B$. Then I simply combine those two finite collections. Their union is still finite, and together they cover both $A$ and $B$, hence they cover $A\cup B$.

So the proof template is: begin with an arbitrary open cover of $A\cup B$, apply compactness once to $A$, apply compactness again to $B$, and then merge the two finite subcovers into one finite subcover of the union.
\end{scratch}

\textbf{Proof.}\\
Let $\{U_i\}_{i\in I}$ be an open cover of $A\cup B$. Thus each $U_i$ is open in $X$ and
\[
A\cup B\subseteq \bigcup_{i\in I} U_i.
\]
Since $A\subseteq A\cup B$, it follows that
\[
A\subseteq \bigcup_{i\in I} U_i.
\]
So $\{U_i\}_{i\in I}$ is an open cover of $A$. Because $A$ is compact, there exists a finite set $J_A\subseteq I$ such that
\[
A\subseteq \bigcup_{i\in J_A} U_i.
\]

Similarly, since $B\subseteq A\cup B$, we also have
\[
B\subseteq \bigcup_{i\in I} U_i.
\]
So $\{U_i\}_{i\in I}$ is an open cover of $B$. Because $B$ is compact, there exists a finite set $J_B\subseteq I$ such that
\[
B\subseteq \bigcup_{i\in J_B} U_i.
\]

Now set
\[
J:=J_A\cup J_B.
\]
Since $J_A$ and $J_B$ are finite, the set $J$ is finite. Moreover,
\[
A\cup B
\subseteq
\left(\bigcup_{i\in J_A} U_i\right)\cup \left(\bigcup_{i\in J_B} U_i\right)
=
\bigcup_{i\in J} U_i.
\]
Thus $\{U_i\}_{i\in J}$ is a finite subcover of $A\cup B$.

Since the open cover $\{U_i\}_{i\in I}$ was arbitrary, we conclude that $A\cup B$ is compact.
\qed


\newpage

\problem[]{3}
Let $(X,d)$ be a metric space.
\begin{enumerate}[label=\textbf{(\alph*)}]
  \item Prove that if $A\subseteq B\subseteq X$, $A$ is closed in $X$, and $B$
  is compact, then $A$ is compact.

  (Hint: Given an open cover of $A$, we can add $A^c$ to get an open cover of
  $B$.)
  \item Prove that if $C,D\subseteq X$ are both compact, then $C\cap D$ is
  compact. (Hint: Use (a).)
\end{enumerate}

\begin{scratch}
Both parts of this problem seem like applications of the open-cover definition of compactness, together with one of the closure results from lecture. For part (a), we are given that $A\subseteq B\subseteq X$, that $A$ is closed in $X$, and that $B$ is compact. Since compactness is about open covers, I should begin with an arbitrary open cover of $A$ and try to turn it into an open cover of $B$. The hint tells me exactly how: because $A$ is closed, its complement $A^c$ is open, so if I add $A^c$ to my cover of $A$, then every point of $B$ is covered either because it lies in $A$ or because it lies in $B\setminus A\subseteq A^c$. Then compactness of $B$ gives a finite subcover, and when I restrict back to $A$, the set $A^c$ becomes unnecessary.

For part (b), I want to prove that the intersection of two compact sets is compact. The cleanest way to use part (a) is to observe that if $C$ and $D$ are compact, then by the theorem from lecture each compact set is closed. In particular, $C$ is closed in $X$, and of course $C\cap D\subseteq D$. Since $D$ is compact, part (a) applies with $A=C\cap D$ and $B=D$. That immediately gives compactness of $C\cap D$.
\end{scratch}

\textbf{Proof.}\\
\textbf{(a)} Let $\{U_i\}_{i\in I}$ be an arbitrary open cover of $A$, where each $U_i$ is open in $X$ and
\[
A\subseteq \bigcup_{i\in I} U_i.
\]
Since $A$ is closed in $X$, the complement $A^c$ is open in $X$. I claim that
\[
\{U_i\}_{i\in I}\cup\{A^c\}
\]
is an open cover of $B$.

Indeed, let $x\in B$. If $x\in A$, then since $\{U_i\}_{i\in I}$ covers $A$, we have $x\in U_i$ for some $i\in I$. If $x\notin A$, then $x\in A^c$. In either case, $x$ lies in one of the sets from $\{U_i\}_{i\in I}\cup\{A^c\}$. Hence
\[
B\subseteq \left(\bigcup_{i\in I} U_i\right)\cup A^c.
\]
So this is an open cover of $B$.

Because $B$ is compact, there exist finitely many indices $i_1,\dots,i_n\in I$ such that
\[
B\subseteq U_{i_1}\cup\cdots\cup U_{i_n}\cup A^c.
\]
Now intersect both sides with $A$. Since $A\subseteq B$, we obtain
\[
A\subseteq A\cap\left(U_{i_1}\cup\cdots\cup U_{i_n}\cup A^c\right).
\]
Distributing the intersection gives
\[
A\subseteq (A\cap U_{i_1})\cup\cdots\cup(A\cap U_{i_n})\cup(A\cap A^c).
\]
But $A\cap A^c=\varnothing$, so
\[
A\subseteq (A\cap U_{i_1})\cup\cdots\cup(A\cap U_{i_n}).
\]
Since each $A\cap U_{i_k}\subseteq U_{i_k}$, it follows that
\[
A\subseteq U_{i_1}\cup\cdots\cup U_{i_n}.
\]
Thus $\{U_{i_1},\dots,U_{i_n}\}$ is a finite subcover of $A$. Since the original open cover of $A$ was arbitrary, we conclude that $A$ is compact.

\textbf{(b)} Suppose that $C,D\subseteq X$ are both compact. By the theorem from lecture, every compact subset of a metric space is closed. Therefore $C$ is closed in $X$.

Now note that
\[
C\cap D\subseteq D\subseteq X.
\]
Since $C$ is closed in $X$, the set $C\cap D$ is also closed in $D$; equivalently, we may simply apply part (a) with
\[
A=C\cap D,
\qquad
B=D.
\]
Indeed, $A\subseteq B$, the set $A=C\cap D$ is closed in $X$ as the intersection of two closed sets, and $B=D$ is compact. Therefore, by part (a), the set $C\cap D$ is compact.
\qed

\newpage

\problem[]{4}
Let $(X,d)$ be a metric space and let $Y\subseteq X$. Prove that if $Y$ is
compact, then $Y$ is complete.

\begin{scratch}
The main definition here is the definition of completeness: a metric space is complete if every Cauchy sequence converges in that space. Since the problem assumes that $Y\subseteq X$ is compact, the idea is to start with an arbitrary Cauchy sequence in $Y$ and show that it must converge to a point of $Y$.

The theorem from lecture that does the heavy lifting is the sequential compactness result: if a set is compact, then every sequence in that set has a subsequence that converges in the set. We also know from lecture that every convergent sequence is Cauchy, and in any metric space a Cauchy sequence can have at most one possible limit.

So the proof template is very standard. I take a Cauchy sequence $\{y_n\}$ in $Y$. Because $Y$ is compact, the sequence has a subsequence $\{y_{n_k}\}$ that converges to some point $y\in Y$. Then I use the fact that the whole sequence is Cauchy to show that actually the entire sequence must converge to that same point $y$. That will prove that every Cauchy sequence in $Y$ converges in $Y$, which is exactly what it means for $Y$ to be complete.
\end{scratch}

\textbf{Proof.}\\
To prove that $Y$ is complete, let $\{y_n\}_{n\ge 1}$ be a Cauchy sequence in the subspace $(Y,d_Y)$. Since the induced metric $d_Y$ is just the restriction of $d$ to $Y$, we may regard $\{y_n\}$ simply as a Cauchy sequence of points in $Y$.

Because $Y$ is compact, every sequence in $Y$ has a subsequence that converges to a point of $Y$. Therefore there exists a subsequence $\{y_{n_k}\}_{k\ge 1}$ and a point $y\in Y$ such that
\[
y_{n_k}\to y.
\]
We claim that in fact the whole sequence $\{y_n\}$ converges to $y$.

Let $\varepsilon>0$. Since $\{y_n\}$ is Cauchy, there exists $N_1\in\mathbb{N}$ such that
\[
d(y_n,y_m)<\frac{\varepsilon}{2}
\qquad \text{for all } n,m\ge N_1.
\]
Since $y_{n_k}\to y$, there exists $N_2\in\mathbb{N}$ such that
\[
d(y_{n_k},y)<\frac{\varepsilon}{2}
\qquad \text{for all } k\ge N_2.
\]
Choose $k\ge N_2$ large enough so that also $n_k\ge N_1$. Now let $n\ge N_1$. By the triangle inequality,
\[
d(y_n,y)
\le d(y_n,y_{n_k})+d(y_{n_k},y).
\]
Since $n\ge N_1$ and $n_k\ge N_1$, the Cauchy condition gives
\[
d(y_n,y_{n_k})<\frac{\varepsilon}{2},
\]
and since $k\ge N_2$, the convergence of the subsequence gives
\[
d(y_{n_k},y)<\frac{\varepsilon}{2}.
\]
Therefore
\[
d(y_n,y)<\frac{\varepsilon}{2}+\frac{\varepsilon}{2}=\varepsilon.
\]
Since this holds for all $n\ge N_1$, we conclude that
\[
y_n\to y.
\]
Because $y\in Y$, the sequence $\{y_n\}$ converges in $Y$.

Thus every Cauchy sequence in $Y$ converges to a point of $Y$, so $Y$ is complete.
\qed


\newpage
% \appendix
% \section*{Appendix: Toolbox}
% \addcontentsline{toc}{section}{Appendix: Toolbox} % I am tired and do not want to do this; it is of little value to me I am realizing

% Example (using amsthm environments instead of boxed tcolorboxes):
% \begin{definition}
% A metric space is ...
% \end{definition}
%
% \begin{theorem}[Heine--Borel]
% ...
% \end{theorem}

\end{document}
